%%%%%%%%%%%%%%%%%%%%%%%%%%%%%%%%%%%%%%%%%%%%%%%%%%%%%%%%%%%%%%%%%%%%%%%%
\section{A branching random walk model}\label{sec:brw}
%%%%%%%%%%%%%%%%%%%%%%%%%%%%%%%%%%%%%%%%%%%%%%%%%%%%%%%%%%%%%%%%%%%%%%%%

A simplified model shows the mechanism of Section~\ref{sec:heuristic} in its purest form. Replace
the dyadic squares by the vertices of the rooted tree in which every vertex has four children. Attach to each vertex
other than the root an independent normal variable with mean
$-Q\log2$ and variance $\log2$, and let the \emph{position} of a vertex be $y_0$ plus the sum of these
variables along the path from the root, where $y_0=Q\log(1/\varepsilon)\ge0$ is the starting
position. A vertex is \emph{alive} if it and all its ancestors have nonnegative positions. Let $V$
be the number of alive vertices; the number of terminal vertices, the children of alive vertices
that are not alive, is $3V+1$.

This is the subdivision of Section~\ref{sec:model}, with the height as the position, except that
the field increments are replaced by independent normal steps; in the real model nearby steps are
correlated. It is a \emph{branching random walk killed at zero}: individuals that fall below zero
are removed together with their descendants, and $V$ is the total number of individuals that ever
live. For one step, the logarithm of four times its exponential moment is
\[
\log4-\theta Q\log2+\frac{\theta^2}2\log2=\psi(\theta)\log2,
\]
with $\psi$ the quadratic~\eqref{eq:quadratic}. For $Q>2$ it has the two roots
$\lambda_\pm=Q\pm\sqrt{Q^2-4}$, that is $\gamma$ and $4/\gamma$.

\begin{theorem}[A\"id\'ekon, Hu and Zindy~\cite{AHZ}]\label{thm:ahz}
For $Q>2$ and the start $y_0=0$, $\Prob(V\ge n)$ is asymptotic to a positive constant times
$n^{-x}$, where $x=\lambda_+/\lambda_-$.
% UNVERIFIED: the project checked the hypotheses of AHZ Theorem 1(ii) for the start 0 only;
% the general-start form of their theorem was not rechecked for this paper.
\end{theorem}

The roots $\lambda_+$ and $\lambda_-$ are the cost and yield rates of Section~\ref{sec:heuristic}. Since $\lambda_+\lambda_-=4$ and $\lambda_++\lambda_-=2Q$,
\[
x+\frac1x=\frac{\lambda_+^2+\lambda_-^2}{\lambda_+\lambda_-}=\frac{(2Q)^2-8}{4}=Q^2-2=\frac{13-c}6 .
\]
So the KPZ quadratic relates the two exponential rates of a killed branching random walk, and
$x=4/\gamma^2$ as in Theorem~\ref{thm:tail}. Here the probability of each size is also controlled.

\begin{theorem}\label{thm:brwlocal}
Let $Q>2$ and $y_0\ge0$. For all $n\ge1$, $\Prob(V=n)$ lies between two positive constant multiples of
$n^{-1-x}$, the constants depending on $Q$ and $y_0$. At $Q^2=29/6$ the exponent is $1+x=\astar$.
\end{theorem}

The lower bound and upper bounds $n^{-1-p}$ for every $p<x$ are proved directly; the sharp upper
bound also uses Theorem~\ref{thm:ahz}. Three ideas make the proof work. First, regard $V=V(y_0)$ as a
nondecreasing step function of the start $y_0$, a \emph{staircase}. Its successive jumps are exactly
independent, each with the law of $V(0)$, because the cluster that a newly alive vertex brings
along depends only on the unused steps below it. So the values of the staircase form a renewal
sequence, and by the renewal theorem each integer is a value with probability tending to
$1/\E[V(0)]>0$: a tail estimate becomes a local one at the cost of a constant. Second, given the
alive set, the margins by which its terminal children miss being alive are independent truncated
normal variables; this controls how long the staircase stays at a value. Third, a tilt along one
path at rate $\lambda_+$ makes one vertex high. The Gaussian free field lacks the independence in
the first two ideas, one reason why the local lower bound for the APPS model is open.
